%
% search-panel.tex
%
% Z Specification for Search Panel UI Coordination
% Formal model of the interaction between search input, results list,
% and detail view in quarry-menubar.
%

\documentclass[a4paper,10pt,fleqn]{article}
\usepackage[margin=1in]{geometry}
\usepackage{fuzz}
\usepackage[colorlinks=true,linkcolor=blue,citecolor=blue,urlcolor=blue]{hyperref}

\begin{document}

\title{Search Panel: A Z Specification}
\author{Formal Model of Search UI Coordination}
\date{March 2026}
\hypersetup{
  pdftitle={Search Panel: A Z Specification},
  pdfauthor={Punt Labs},
  pdfsubject={Z Specification},
  pdfcreator={fuzz/probcli}
}
\maketitle

\tableofcontents
\newpage

\section{Introduction}

This specification models the coordination between three UI components in
\texttt{quarry-menubar}: the search input field, the results list, and the
detail view. The content pane shows either the results list or a single
result's detail --- never both.

The key property is that clearing the search input must reset all
downstream state: results, the selected result (detail view), and the
highlighted result (list cursor). The current implementation violates this
invariant when the query is already empty and a programmatic clear is
issued, because the reactive change handler only fires on value
\emph{transitions}, not on assignments to the same value.

\section{Basic Types}

\begin{zed}
  [RESULTID]
\end{zed}

A result identifier. In the implementation, each search result carries a
unique ID used for list selection and detail navigation.

\section{Free Types}

\begin{zed}
  SearchState ::= idle | loading | hasResults | noResults | searchError
\end{zed}

The lifecycle of a search request. \texttt{idle} means no query is active.
\texttt{loading} means a request is in flight. The remaining three are
terminal states for a completed search.

\begin{zed}
  ZBOOL ::= ztrue | zfalse
\end{zed}

\section{Global Constants}

\begin{axdef}
  maxResults : \nat
  \where
  maxResults = 50
\end{axdef}

Upper bound on the number of results returned by a single search. Exists
for ProB animation --- the actual implementation has no hard cap, but
the semantic search API returns a bounded result set.

\section{State}

A single flat schema captures all UI coordination state. The search
input, search lifecycle, result set, detail selection, and list cursor
are modeled together so that their invariants are expressed in one place.

\begin{schema}{State}
  queryEmpty : ZBOOL \\
  searchState : SearchState \\
  results : \finset RESULTID \\
  selectedResult : \finset RESULTID \\
  highlightedResult : \finset RESULTID
  \where
  %
  % Optional values: at most one element
  %
  \# selectedResult \leq 1 \\
  \# highlightedResult \leq 1 \\
  %
  % Selection must come from the result set
  %
  selectedResult \subseteq results \\
  highlightedResult \subseteq results \\
  %
  % Results only exist when search has completed with results
  %
  searchState \neq hasResults \implies results = \emptyset \\
  %
  % CRITICAL INVARIANT: empty query clears everything downstream
  %
  queryEmpty = ztrue \implies searchState = idle \\
  queryEmpty = ztrue \implies selectedResult = \emptyset \\
  queryEmpty = ztrue \implies highlightedResult = \emptyset \\
  %
  % Cannot select a result for detail unless results exist
  %
  searchState \neq hasResults \implies selectedResult = \emptyset \\
  %
  % Detail view and list cursor are mutually exclusive:
  % when viewing detail, the list is not rendered, so no highlight
  %
  selectedResult \neq \emptyset \implies highlightedResult = \emptyset \\
  %
  % Bounded result set for animation
  %
  \# results \leq maxResults
\end{schema}

\begin{schema}{Init}
  State'
  \where
  queryEmpty' = ztrue \\
  searchState' = idle \\
  results' = \emptyset \\
  selectedResult' = \emptyset \\
  highlightedResult' = \emptyset
\end{schema}

\section{Operations}

\subsection{EnterQuery}

The user types a non-empty string into the search field. This starts a
debounced search. The search state moves to \texttt{loading} and any
previous results and selections are cleared.

\begin{schema}{EnterQuery}
  \Delta State
  \where
  queryEmpty' = zfalse \\
  searchState' = loading \\
  results' = \emptyset \\
  selectedResult' = \emptyset \\
  highlightedResult' = \emptyset
\end{schema}

\subsection{ClearQuery}

The user deletes all text from the search field (backspace, select-all
then delete, or any other mechanism). This is the operation that must
enforce the critical invariant: everything downstream resets.

\begin{schema}{ClearQuery}
  \Delta State \\
  \where
  queryEmpty = zfalse \\
  queryEmpty' = ztrue \\
  searchState' = idle \\
  results' = \emptyset \\
  selectedResult' = \emptyset \\
  highlightedResult' = \emptyset
\end{schema}

\subsection{ReceiveResults}

An asynchronous search completes with a non-empty result set. The query
must still be non-empty (the user has not cleared it while the search was
in flight).

\begin{schema}{ReceiveResults}
  \Delta State \\
  newResults? : \finset RESULTID
  \where
  searchState = loading \\
  queryEmpty = zfalse \\
  newResults? \neq \emptyset \\
  \# newResults? \leq maxResults \\
  searchState' = hasResults \\
  results' = newResults? \\
  selectedResult' = \emptyset \\
  highlightedResult' = \emptyset \\
  queryEmpty' = queryEmpty
\end{schema}

\subsection{ReceiveEmpty}

A search completes with no results.

\begin{schema}{ReceiveEmpty}
  \Delta State
  \where
  searchState = loading \\
  queryEmpty = zfalse \\
  searchState' = noResults \\
  results' = \emptyset \\
  selectedResult' = \emptyset \\
  highlightedResult' = \emptyset \\
  queryEmpty' = queryEmpty
\end{schema}

\subsection{ReceiveError}

A search fails.

\begin{schema}{ReceiveError}
  \Delta State
  \where
  searchState = loading \\
  queryEmpty = zfalse \\
  searchState' = searchError \\
  results' = \emptyset \\
  selectedResult' = \emptyset \\
  highlightedResult' = \emptyset \\
  queryEmpty' = queryEmpty
\end{schema}

\subsection{SelectResult}

The user clicks a result (or presses Enter on a highlighted result) to
open the detail view. The result must be in the current result set.
Highlighting is cleared because the detail view replaces the list.

\begin{schema}{SelectResult}
  \Delta State \\
  result? : RESULTID
  \where
  searchState = hasResults \\
  result? \in results \\
  selectedResult' = \{ result? \} \\
  highlightedResult' = \emptyset \\
  results' = results \\
  searchState' = searchState \\
  queryEmpty' = queryEmpty
\end{schema}

\subsection{HighlightResult}

The user arrow-keys to highlight a result in the list. Only available
when the detail view is not showing (selectedResult is empty).

\begin{schema}{HighlightResult}
  \Delta State \\
  result? : RESULTID
  \where
  searchState = hasResults \\
  selectedResult = \emptyset \\
  result? \in results \\
  highlightedResult' = \{ result? \} \\
  selectedResult' = selectedResult \\
  results' = results \\
  searchState' = searchState \\
  queryEmpty' = queryEmpty
\end{schema}

\subsection{ClearHighlight}

The user presses Escape when a result is highlighted but no detail is
open. The highlight is cleared.

\begin{schema}{ClearHighlight}
  \Delta State
  \where
  selectedResult = \emptyset \\
  highlightedResult \neq \emptyset \\
  highlightedResult' = \emptyset \\
  selectedResult' = selectedResult \\
  results' = results \\
  searchState' = searchState \\
  queryEmpty' = queryEmpty
\end{schema}

\subsection{CloseDetail}

The user presses Escape or the Back button while viewing a detail. The
detail view closes and the results list is shown again. The query and
results are preserved.

\begin{schema}{CloseDetail}
  \Delta State
  \where
  selectedResult \neq \emptyset \\
  selectedResult' = \emptyset \\
  highlightedResult' = \emptyset \\
  results' = results \\
  searchState' = searchState \\
  queryEmpty' = queryEmpty
\end{schema}

\subsection{ProgrammaticClear}

An external trigger clears the search state (e.g., switching databases).
This must reset everything regardless of the current query value --- this
is where the implementation bug manifests, because the reactive handler
only fires on value \emph{changes}.

\begin{schema}{ProgrammaticClear}
  \Delta State
  \where
  queryEmpty' = ztrue \\
  searchState' = idle \\
  results' = \emptyset \\
  selectedResult' = \emptyset \\
  highlightedResult' = \emptyset
\end{schema}

\subsection{ChangeCollection}

The user picks a different collection filter. This clears selection and
restarts the search if a query exists.

\begin{schema}{ChangeCollection}
  \Delta State
  \where
  selectedResult' = \emptyset \\
  highlightedResult' = \emptyset \\
  searchState' = \IF queryEmpty = ztrue \THEN idle \ELSE loading \\
  results' = \emptyset \\
  queryEmpty' = queryEmpty
\end{schema}

\section{System Invariants}

The specification enforces these properties across all reachable states:

\begin{enumerate}
  \item \textbf{Empty query implies idle}: no search is in progress and no
    results are displayed when the input field is empty.
  \item \textbf{Empty query implies no selection}: neither the detail view
    nor the list cursor can be active when the input field is empty. This
    is the invariant violated by the current implementation.
  \item \textbf{Selection comes from results}: you cannot view the detail
    of a result that is not in the current result set.
  \item \textbf{Detail and highlight are exclusive}: the detail view
    replaces the results list, so both cannot be active simultaneously.
  \item \textbf{Results require hasResults state}: the result set is
    non-empty only when the search state is \texttt{hasResults}.
\end{enumerate}

\end{document}
